% ===== tr-07 块1：第 5 章 pp.1--4（原文页 1--4）=====

% [原文 p.1]
\section*{摘要}

本文综述 Strominger--Yau--Zaslow 猜想在度量方面的若干侧面，该猜想涉及 Calabi--Yau 流形在大复结构极限附近特殊拉格朗日纤维化的存在性。我们将讨论这一猜想图景的种种动机、最好的期望是什么，以及若干微妙之处。综述的主体部分着重说明多复变势理论、尤其是非阿基米德几何所起的作用，并列出一系列公开问题。

\section{概述}

本文综述 Strominger--Yau--Zaslow 猜想在度量方面的近期进展，该猜想关乎 Calabi--Yau 流形在大复结构极限附近特殊拉格朗日纤维化的存在性。本文主要基于 [52][53][51] 以及随后的若干报告。第 5.2--5.6 节的内容大略如下：

\begin{itemize}
\item 我们从镜像对称、极小曲面、黎曼几何、复几何与非阿基米德几何出发，追溯 SYZ 猜想的多种动机。我们先讨论最乐观的解读及其困难，随后转向更为谨慎的弱版本。
\item 我们从复几何的角度回顾大复结构极限与本质骨架的含义，随后简要综述 Kontsevich--Soibelman 猜想。
\item 我们解释各种分析工具。文中简要概述 Yau 对 Calabi 猜想的解决，以说明为何理解大复结构极限需要新的想法。我们的主要着眼点是复势理论（complex pluripotential theory）------它是 [52][53] 的分析核心------此外还有椭圆偏微分方程中 Savin 的小扰动定理，以及实 Monge--Ampère 方程的正则性理论。
% [原文 p.2]
\item 我们以最精简的方式概述非阿基米德势理论：Berkovich 空间、半正度量（semipositive metrics）以及非阿基米德 Calabi 度量的概念；并解释猜想中的“比较性质”（comparison property）背后的直觉，这一性质是赋予非阿基米德 Calabi 度量以微分几何诠释所必需的。
\item 我们勾勒近期进展 [52][53] 的证明，着重于其中的直觉、微妙之处与公开问题。
\end{itemize}

% [Tutor Note: 原文 p.1 第 1 节首句把 Strominger 误拼作 “Stominger”（600 dpi 裁图核实：p.1，x 100--520、y 405--440 PDF 点）；译文按正确拼写 Strominger--Yau--Zaslow 处理，不改原文信息。]

\para{致谢。}作者现任 Clay Research Fellow 与 MIT CLE Moore Instructor。他感谢 Léonard Pille-Schneider 的评论。

\section{度量 SYZ 猜想}

Strominger--Yau--Zaslow 猜想最初由物理与微分几何两方面的考量共同驱动，它处在镜像对称、极小曲面理论、黎曼几何、复 Kähler 几何与非阿基米德几何的十字路口。我们将追溯若干在历史上有重要意义的进展，看看对分析困难的逐渐认识如何导致对猜想作出各种折中的解读。

\subsection{SYZ 猜想的起源}

一个 $n$ 维 Calabi--Yau（CY）流形 $(X,\omega,\Omega)$ 是指带有处处非零全纯体积形式 $\Omega$、并满足复 Monge--Ampère（MA）方程
\begin{equation}
\omega^n = \mathrm{const}\,\Omega\wedge\overline{\Omega} \tag{5.1}
\end{equation}
的 Kähler 流形；该方程蕴含度量的 Ricci 平坦性。此外，这类流形容许平行旋量，因而可以作为超对称 II 型弦理论中目标空间度量的候选。相位为 $\theta$ 的特殊拉格朗日子流形是满足
\begin{equation}
\omega|_L = 0,\qquad \mathrm{Im}(e^{-i\theta}\Omega)|_L = 0 \tag{5.2}
\end{equation}
的 $n$ 维子流形 $L$。由于校准不等式
\begin{equation}
\int_L \mathrm{Re}(e^{-i\theta}\Omega) \le \int_L dvol_L = \mathrm{Vol}(L) \tag{5.3}
\end{equation}
恰好被特殊拉格朗日取到等号，它们在各自的同调类中是绝对极小化子。在物理上，它们对应于 BPS D-膜的支撑。

% [Tutor Note: 式 (5.1) 中的 const 依赖于 $\Omega$ 的归一化：取定 $\Omega$ 后该式即 $\omega^n=c\,\Omega\wedge\overline{\Omega}$，“Ricci 平坦”这一等价性要求 $c$ 与 $\Omega$ 的选取相容；原文用 const 而未写出具体常数，是因为后文对体积另有归一化约定。]

% [Tutor Note: 式 (5.3) 的第一个不等号来自校准形式 $\mathrm{Re}(e^{-i\theta}\Omega)$ 在 $L$ 上不超过体积元；等号成立当且仅当 $L$ 为相位 $\theta$ 的特殊拉格朗日，这正是原文 saturated precisely by 的含义。相位只能取模 $2\pi$。]

\textbf{Strominger--Yau--Zaslow 猜想} [67] 的原始形式是问：

% [原文 p.3]
\begin{conjecture}{5.2.1}
给定紧 Calabi--Yau 流形 $X$，当它接近大复结构极限时，能否在 $X$ 上找到特殊拉格朗日环面纤维化？
\end{conjecture}

SYZ 猜想 [67] 的物理起源在很大程度上来自镜像对称，其大略梗概如下。由同调镜像对称，人们期望紧 Calabi--Yau 流形 $X$ 有一个镜像 $X^\vee$，使得两边的 D-膜范畴被等同起来。在全纯侧 $X^\vee$（“B-侧”），点 $x\in X^\vee$ 支撑摩天楼层 $\mathcal{O}_x$，它们应当对应于辛侧 $X$（“A-侧”）内部的某些拉格朗日膜。扩张群 $\operatorname{Ext}^*(\mathcal{O}_x,\mathcal{O}_x)\simeq H^*(T^n)$，这提示这些拉格朗日膜是环面对象。当 $x\ne y$ 时，$\mathcal{O}_x$ 与 $\mathcal{O}_y$ 之间的 Ext 群应当消失，这（并不具有决定性\footnote{Floer 上同调的消失并不蕴含 Floer 上链空间的消失，而且似乎没有强有力的论证能排除特殊拉格朗日相交。}）提示这些环面互不相交，从而引出拉格朗日纤维化结构的猜测。然而关于“特殊”拉格朗日的论断超出了纯粹的同调镜像对称，它来自 D-膜上的 BPS 条件。所有 $\mathcal{O}_x$ 的模空间就是镜像流形 $X^\vee$，它随后应当与支撑在特殊拉格朗日环面上的 BPS 膜的模空间等同。这种模空间解释给出镜像流形 $X^\vee$ 上 Kähler 结构的零阶近似，但须受与全纯圆盘（“瞬子修正”）相关的高阶修正的影响，后者的效应除在奇异纤维附近之外被认为按指数衰减。若忽略奇异纤维带来的微妙之处，SYZ 图景便提供了一套重建镜像、并把同调镜像对称解释为某种 Fourier--Mukai 变换（“镜像对称即 T 对偶”）的方案。

% [Tutor Note: 原文 p.3 把 skyscraper sheaf 拼作 “skyscrapper sheaves”，拼写有误（600 dpi 裁图核实同页）；译文按标准术语“摩天楼层（skyscraper sheaf）”处理。]

\para{记号。}我们的约定是 $d=\partial+\bar{\partial}$，$d^c=\frac{\sqrt{-1}}{2\pi}(-\partial+\bar{\partial})$，于是 $dd^c=\frac{\sqrt{-1}}{\pi}\partial\bar{\partial}$。Kähler 势与 Kähler 度量之间的关系是 $\omega_\phi=\omega+dd^c\phi$。另一种做法是用局部绝对势来理解 Kähler 度量，即取 $\omega=dd^c\varphi$，其中 $\varphi$ 是局部定义的多重次调和（psh）函数。给定线丛 $L$ 上的 Hermitian 度量 $h$，其曲率形式为 $-dd^c\log h^{1/2}$，位于类 $c_1(L)$ 中。

% [Tutor Note: 原文的 $d^c$ 归一化使 $dd^c=\frac{\sqrt{-1}}{\pi}\partial\bar{\partial}$，恰为常见约定 $\frac{\sqrt{-1}}{2\pi}\partial\bar{\partial}$ 的两倍；正因如此，Hermitian 度量 $h$ 的曲率形式才写成 $-dd^c\log h^{1/2}$ 而非常见的 $-dd^c\log h$。两处 $1/2$ 与 $d^c$ 归一化相互抵消，并非笔误（600 dpi 裁图核实：p.3，y 440--480 PDF 点）。]

从微分几何上看，SYZ 原文所给出的主要证据是\textbf{半平坦度量}。考虑对数映射 $\operatorname{Log}_t:(\C^*)^n\to\R^n$，其定义域取在某个开凸子集 $U\subset\R^n$ 之上，
\[
\operatorname{Log}_t(z_1,\dots z_n) = \frac{1}{\log|t|}(\log|z_1|,\dots\log|z_n|).
\]
在加上 $T^n$ 对称性之后，由初等的 Hessian 计算可知：$\phi$ 在底空间 $U$ 上是光滑严格凸函数，当且仅当它拉回到 $\operatorname{Log}_t^{-1}(U)$ 上是光滑 Kähler 势；并且 Calabi--Yau 条件
\[
(dd^c\phi\circ\operatorname{Log}_t)^n = \frac{\mathrm{const}}{|\log|t||^{2n}}\prod_i d\log z_i\wedge d\overline{\log z_i}
\]
% [原文 p.4]
等价于实 Monge--Ampère 方程
\[
\det(D^2\phi)=\mathrm{const}.
\]

% [Tutor Note: 上式右端的因子 $|\log|t||^{-2n}$ 与 $\operatorname{Log}_t$ 的分母 $\log|t|$ 配套：$dd^c$ 作用在只依赖 $\log|z_i|/\log|t|$ 的函数上会带出两个 $1/\log|t|$，$n$ 次幂正好给出 $2n$ 次。抄写时易把分母误作 $|\log|t||^n$ 或漏掉绝对值中的绝对值（外层 $|\cdot|$ 对内层 $\log|t|$ 再取绝对值）。]

在此情形下，度量 $dd^c\phi\circ\operatorname{Log}_t$ 称为半平坦的，这是因为 $T^n$ 纤维上的限制由于 $T^n$ 对称性而是欧氏的。相对于 Calabi--Yau 结构
\[
\omega = dd^c(\phi\circ\operatorname{Log}_t),\qquad \Omega=\sqrt{-1}^{\,n}\prod d\log z_i,
\]
这些 $T^n$ 纤维是相位零的特殊拉格朗日。引入归一化参数 $t$ 的目的是：当 $t\to0$ 时，$T^n$ 纤维收缩到零尺寸，而 Calabi--Yau 度量 $dd^c(\phi\circ\operatorname{Log}_t)$ 收敛到 $U\subset\R^n$ 上的实 Monge--Ampère 度量
\begin{equation}
g_0 = \frac{1}{2\pi}\sum_{i,j}\frac{\partial^2\phi}{\partial x_i\partial x_j}dx_idx_j,\qquad \det(D^2\phi)=\mathrm{const}. \tag{5.4}
\end{equation}

% [Tutor Note: 式 (5.4) 中的 $\Omega=\sqrt{-1}^{\,n}\prod d\log z_i$ 是 $(\C^*)^n$ 上的全纯体积形式，其指数是 $n$；不要与 §3 中 $\Omega_t\wedge\overline{\Omega_t}$ 前的归一化因子 $\sqrt{-1}^{\,n^2}$（指数为 $n^2$）混淆。二者来源不同，抄写时须逐处对照页图。]

现在考虑大复结构极限附近的情形。大复结构极限是 Calabi--Yau 度量族的一类特定极限情形，人们期望半平坦度量在 Calabi--Yau 流形的一般区域中作为退化 Calabi--Yau 度量的渐近描述出现。在 SYZ 图景中，一般区域按启发式理解就是指远离奇异特殊拉格朗日纤维的区域。要点在于：在一般区域中，特殊拉格朗日只是局部环面坐标卡中对数映射的小扰动。当我们趋近大复结构极限时，一般区域所占 Calabi--Yau 体积测度的百分比应当趋于 $100\%$。

% [Tutor Note: 这一段关于“一般区域占比趋于 $100\%$”的陈述是启发式的期望，原文并未给出其证明或精确定义；“generic region” 在 SYZ 文献中常译“一般区域”，与 algebraic geometry 中的 “general” 不同。]

这段对 SYZ 启发式推理的复述强调了几点需要谨慎之处：

\begin{itemize}
\item 度量方面的预测在\emph{一般区域}中更有说服力。非一般区域受瞬子修正效应的影响，其度量意义要难分析得多。
\item SYZ 纤维化很可能是大复结构极限附近才\emph{涌现}的行为。特别地，该猜想关心的是 Calabi--Yau 流形的一个\emph{族}，诸如半平坦度量这样的特征只会在足够接近极限时才出现。
\end{itemize}

\subsection{进一步的动机}

除了在镜像对称中的重要性之外，SYZ 猜想对极小曲面理论、黎曼几何、Kähler 几何与代数几何等其他多个领域也很有意义。

\begin{itemize}
\item 特殊拉格朗日是极小子流形，事实上还是校准子流形。目前，在足够多的 Calabi--Yau 流形上构造特殊拉格朗日的方法还很少，尽管存在由 Thomas--Yau [75][74] 开创、并在 [43][57] 中进一步发展的系列猜测。
\end{itemize}
